The US is looking into the deployment of advanced reactors that 
require uranium enriched betwee 5-20\% uranium-235, often referred to 
as \gls{HALEU}. There are no facilities 
in the US to produce \gls{HALEU} in the front end of the fuel cycle, 
prompting questions of how \gls{HALEU} will be supplied. 
\gls{HALEU} can be produce through two primary methods: downblending 
\gls{HEU} and enriching natural uranium. Downblending \gls{HEU} is 
limited by the amount of \gls{HEU} and the impurities present in the 
\gls{HEU}. Enriching natural uranium is limited by the \gls{SWU} 
capacity to enrich the uranium and the amount of natural uranium 
available. Therefore, to understand how each of these production 
methods can be used and the facilities need for each method the 
material requriements of transitioning to \gls{HALEU}-fueled 
reactors can be quantified. 

The purpose of this prelim is to investigate the effects of deploying 
\gls{HALEU}-fueled reactors in the US. This includes quantifying 
the material requirements of potential transitions to these 
reactors. The work performed so far models the transition from the
current fleet of \glspl{LWR} to different subsets of advanced reactors
in a once-through fuel cycle. 
\gls{HALEU}-fueled reactors considered include the X-energy Xe-100 and 
the \gls{USNC} \gls{MMR}. The NuScale VOYGR is also included to understand 
how the deployment of a non-\gls{HALEU}-fueled reactor in tandum with 
\gls{HALEU}-fueled reactors affects the fuel cycle needs. The 
transition is modeled assuming both a no growth and a 1\% 
annual growth in energy demand. Materials of interest include the 
number of reactors built, the mass of enriched uranium, the mass 
of feed uranium to produce the enriched uranium, the \gls{SWU} 
capacity needed, and the amount of \gls{SNF} that must be disposed of. 

The results from analyzing the once-through transition scenarios show 
that each of the material requirements are most dependent on the 
reactors deployed in the scenario and less sensitive to the energy 
demand of the scenerario. The number of reactors needed in each scenarios 
scales with the power output of the reactors deployed in each scenario 
and the energy demand, 
with the transition to only the \gls{MMR} requiring the most reactors for 
each energy demand modeled. The most resources to produce 
\gls{HALEU} are also required by scenarios that deploy only the \gls{MMR}
because of the large number of reactors needed. Deploying the VOYGR 
alongside the \gls{HALEU}-fueled reactors consistently increases the 
mass of enriched uranium and waste compared with the scenarios with 
only \gls{HALEU}-fueled reactors. The increase is a result of the 
VOYGR requiring more fuel in the core than the other two reactors. 
For the feed uranium and \gls{SWU} capacity needed, there is no consistent 
change when deploying the VOYGR alongside the \gls{HALEU}-fueled reactors. 
The change is dependent on how many VOYGRs are deployed in relation to 
the \gls{HALEU}-fueled reactors. All material metrics for \gls{HALEU} 
(mass of \gls{HALEU}, feed uranium, \gls{SWU} capacity, and \gls{HALEU} 
waste) all decrease when VOYGRs are deployed with \gls{HALEU} reactors
because the VOYGR does not require \gls{HALEU}.

The proposed work for the remainder of this investigation builds upon 
the work already completed to understand the impacts of deploying 
\gls{HALEU}-fueled reactors. The next step is to model the transition
to these reactors in a fuel cycle with recycling. This step will 
quantify the material requirements for deploying these reactors in 
a different fuel cycle option. The next step is to perform sensitivity 
analysis on some of the modeled transitions, and use the results of 
the analysis to identify optimized transitions. Sensitivity anlaysis 
will identify the impact of input parameters of the transitions affect 
the material requirements of the transition. Finally, the effect of 
impurities in \gls{HALEU} created by downblending \gls{HEU} on 
reactor performance will be investigated. This step of the investigation 
will provide insight into if these impurities will detract from expected 
reactor performance enough to warrant limiting the amount of this 
\gls{HALEU} that can be used in a reactor. 